\section{Localization at a root plane}
\label{sec:localization}

We now establish the local description needed for the edge comparison.
Localization along reflection-subgroup cosets appears in Dyer's nil-Hecke
and projective-module approach to Schubert singularities
\cite[Sections~1.7--1.9 and~3.16]{DyerRankTwo}. Here we prove the
precise decomposition for the moment-graph sheaf and its actual edge maps.
The key point is that localization leaves a dihedral moment graph, and on
that graph every indecomposable Braden--MacPherson sheaf is a structure
sheaf.  We will also prove the resulting direct-sum decomposition over the
local coefficient ring; no preservation of indecomposability under
localization is asserted or needed.

Throughout this section, $A$ is the polynomial ring of the real
realization fixed in Lemma~\ref{bg:realization}, with roots in degree
one.  The sheaf $B$ has top $b$ and is supported on $[1,b]$.
Its stalks and edge modules are denoted by $B^a$ and $B^e$, respectively.
We use the generation, flabbiness, and upper-end quotient properties from
Theorem~\ref{bg:bmp}, with the normalization of
Proposition~\ref{bg:normalization}.

\subsection{Dihedral structure sheaves}

We first record the elementary polynomial calculation that will identify
the dihedral models.

\begin{lemma}[Dihedral Kazhdan--Lusztig polynomials]
\label{loc:dihedral-kl}
Let $D$ be a finite or infinite dihedral Coxeter group.  For all $a\le_D c$,
\[
  P^D_{a,c}(q)=1.
\]
\end{lemma}

\begin{proof}
Use the Hecke algebra convention
$(T_s+1)(T_s-q)=0$. The bar involution sends $q^{1/2}$ to $q^{-1/2}$
and $T_s$ to $T_s^{-1}$ for each simple generator, preserving products.
For $w\in D$, put
\[
  C_w^{\circ}=q^{-\ell_D(w)/2}\sum_{z\le_D w}T_z.
\]
We prove that these elements are bar invariant.  This is immediate for
$w=1$ and for a simple reflection.  Let $w$ have an alternating reduced
expression starting with $s$, and set $n=\ell_D(w)\ge2$.
We claim that
\begin{equation}
\label{loc:dihedral-product}
 C_s^{\circ}C_{sw}^{\circ}
 =
 \begin{cases}
  C_w^{\circ},&n=2,\\
  C_w^{\circ}+C_p^{\circ},&n>2,
 \end{cases}
\end{equation}
where $p$ is the prefix of length $n-2$ of the chosen expression for $w$.

To check this, partition $D$ into left $s$-strings $\{r,sr\}$ with
$\ell_D(sr)=\ell_D(r)+1$.  A lower interval cannot contain $sr$ without
containing $r$.  If only $r$ is present in $[1,sw]_D$, then
\[
 (1+T_s)T_r=T_r+T_{sr}.
\]
If both elements are present, the quadratic relation gives
\[
 (1+T_s)(T_r+T_{sr})=(1+q)(T_r+T_{sr}).
\]
The strings that meet $[1,sw]_D$ have union $[1,w]_D$.
There are no doubly included strings when $n=2$.  When $n>2$, their union
is $[1,p]_D$: both alternating words occur in every length below $n-2$,
and the unique element of length $n-2$ in this union is the prefix $p$.
This is also valid when $w$ is the longest element of a finite dihedral
group.  Consequently
\[
 (1+T_s)\sum_{z\le_D sw}T_z
 =\sum_{z\le_D w}T_z+
   \begin{cases}
    0,&n=2,\\
    q\displaystyle\sum_{z\le_D p}T_z,&n>2.
   \end{cases}
\]
Multiplication by $q^{-n/2}$ proves
\eqref{loc:dihedral-product}.  Induction now proves bar invariance of
$C_w^{\circ}$.

Relative to the normalized standard basis
$\{q^{-\ell_D(z)/2}T_z:z\in D\}$, the leading coefficient of
$C_w^{\circ}$ is $1$, and its lower coefficients are
\[
 q^{-(\ell_D(w)-\ell_D(z))/2}
 \in q^{-1/2}\mathbb Z[q^{-1/2}]
 \qquad(z<_D w).
\]
Bar invariance and the triangular characterization therefore give
\[
 C_w^{\circ}
 =q^{-\ell_D(w)/2}\sum_{z\le_D w}P^D_{z,w}(q)T_z.
\]
Comparison of coefficients proves the lemma.
\end{proof}

Let $\Pi$ be a plane spanned by two independent roots, and let $D$ be its
maximal dihedral reflection subgroup.  We always use the canonical
positive roots of $D$ induced from the ambient system and the relative
order on each coset from Lemma~\ref{dih:plane}.
The associated reflection-subgroup description is valid for arbitrary
Coxeter systems by Lemma~\ref{dih:plane}; Dyer's finite-setting treatment
is a predecessor \cite{Dyer91}.
For a coset $C$ and a vertex $c\in C$, define a sheaf
$\mathcal O_c$ on the relative dihedral graph of $C$ by
\[
 \mathcal O_c^a=
 \begin{cases}
 A,&a\le_D c,\\
 0,&a\not\le_D c.
 \end{cases}
\]
On an edge whose two endpoints lie below $c$, its module is
$A/(\alpha_e)$ and both restrictions are the quotient map.  On every
other edge its module is zero.  This is consistent with the upper-end
quotient rule because its support is a relative lower ideal.

\begin{lemma}[Dihedral models in the induced realization]
\label{loc:models}
The Braden--MacPherson sheaf with top $c$ on the relative dihedral graph
of $C$, normalized to have top stalk $A$ in degree zero, is isomorphic to
$\mathcal O_c$.  In particular $\mathcal O_c$ is generated at its stalks by
sections and is flabby for relative upper sets.
The statement uses the actual ambient root labels, including when $D$ is
infinite and its induced representation is not the standard geometric
representation of an infinite dihedral group.
\end{lemma}

\begin{proof}
We first justify the character theorem for precisely this realization.
Restrict the ambient finite-dimensional representation to $D$.
It is faithful.  Every reflection of $D$ has rank-one change from the
identity.  Conversely, suppose an element $d\in D$ has rank-one change.
Ambient reflection faithfulness implies that $d$ is an ambient
reflection.  The image of $d-1$ lies in $\Pi$, since this is true for every
product of reflections whose roots lie in $\Pi$.  Its root therefore
lies in $\Pi$, and the maximal plane-subgroup description makes it a
reflection of $D$.  The restricted representation is consequently
reflection faithful.

If $D$ is finite of type $I_2(m)$, its normalized canonical simple roots
have Gram matrix with diagonal entries one and off-diagonal entries
$-\cos(\pi/m)$ by Lemma~\ref{dih:midpoint-order}. Its root plane is therefore
positive definite and
carries the usual real dihedral reflection representation.  The nondegenerate
ambient space splits into this plane and its orthogonal complement,
and $D$ fixes the latter pointwise.  The real Soergel character theorem
therefore applies to the plane representation \cite{EW14}; adjoining
the fixed polynomial directions preserves the character formula and the
Braden--MacPherson construction.  Indeed the relevant kernels and images
commute with this flat polynomial extension, and graded minimal free
covers remain minimal.

If $D$ is infinite, it is a universal Coxeter system of rank two.
Fiebig's universal-group character theorem applies under its standing
reflection-faithful hypotheses and orbit choices
\cite[Definition~2.1, Sections~2 and~5, and Theorem~9.3]{Fiebig08}.
These requirements hold for the restricted real representation. A regular
orbit exists by avoiding the countable union of proper fixed subspaces.
The required simple-subregular orbits exist by the generic-hyperplane
argument in Lemma~\ref{bg:realization}, applied to the restricted
orthogonal representation.
In particular, no assertion that its simple-root pairing is the
geometric pairing is needed.
The moment-graph identification and
duality give the same stalk normalization as above
\cite[Corollary~6.5 and Proposition~7.1(4)]{Fiebig08}.

The coset identification preserves the relative orientation and the
left root labels, so these statements apply on $C$.
Lemma~\ref{loc:dihedral-kl} and the stalk character formula now show that
every stalk below $c$ is $A$ with a generator in degree zero.
Every nonzero edge module is likewise $A/(\alpha_e)$ with its generator
in degree zero.

Choose a homogeneous global section extending the generator at $c$.
We show that its value at every vertex of $[m,c]_D$ is a generator, where
$m$ is the relative minimum of $C$.  Descend along any path from $c$ to
the chosen vertex.  If its value at the upper endpoint of an edge is a
nonzero degree-zero generator, its restriction to that edge is nonzero,
by the upper-end quotient rule.  Compatibility forces the section value
at the lower endpoint to be nonzero.  Since that stalk is a rank-one
free module generated in degree zero, a nonzero homogeneous value of
degree zero is again a generator.  This proves the assertion along the
path, and every vertex below $c$ lies on such a path.

Use these section values as the generators of the vertex stalks, and
use their common restrictions as the generators of the edge modules.
Both maps at every supported edge are then the quotient map.  This
identifies the sheaf with $\mathcal O_c$.  Generation and flabbiness are
properties of the Braden--MacPherson sheaf, so they hold for this model.
\end{proof}

\subsection{Relative flabbiness after localization}

Let $\mathfrak p\subset A$ be the prime ideal generated by the plane
$\Pi$, and put
\[
 R=A_{\mathfrak p},\qquad
 \mathfrak m=\mathfrak pR,\qquad
 B_R=B\otimes_A R.
\]
The ring $R$ is a regular local unique factorization domain of dimension
two.  We regard its sheaves as ungraded.
An edge module of $B_R$ vanishes unless its label lies in
$\Pi$: any other label becomes a unit.  The surviving edges thus lie
inside the cosets of $D$.

Fix such a coset $C$ and write
\[
 K=C\cap[1,b].
\]
This is a finite relative lower ideal.  Indeed, if $a\le_D c$ with
$c\in K$, every relative upward edge is also an ambient upward edge,
so $a\le c\le b$ in the ambient order and hence $a\in K$.
This uses only the implication from relative to ambient comparability;
the two orders on the coset need not coincide.

\begin{lemma}
\label{loc:relative-flabbiness}
The restriction $B_R|_K$ has free stalks, the upper-end quotient
rule, generation at each stalk by sections, and flabbiness for upper
sets in the relative dihedral order.
\end{lemma}

\begin{proof}
Freeness and the quotient rule persist under localization.  The graph
on $[1,b]$ is finite, and its section module is the kernel of a map
between finite direct sums of stalk and edge modules.  Formation of
sections therefore commutes with flat localization.  In particular,
ambient flabbiness and generation at stalks persist.  Restricting a
global section to $K$ proves generation there.

To prove relative flabbiness, let $J\subseteq K$ be a relative upper
set and prescribe a section on $J$.  Regard $J$ as a subset of $[1,b]$.
It is upper closed under all surviving edges: those edges stay in $C$,
and their orientations agree with the relative dihedral graph.
We extend the section by adding missing vertices in decreasing ambient
order.

When a missing vertex $a$ is added, all ambient vertices strictly above
$a$ are already present.  Put
\[
 \Omega_a=\{z\in[1,b]:a<z\}.
\]
This is an ambient upper set.  Ambient flabbiness extends the restriction
of the prescribed section to $\Omega_a$; use the resulting value at $a$.
It is compatible with every already prescribed upper neighbor.
An already prescribed lower neighbor $z$ cannot be joined to $a$ by a
surviving edge: upper closure under surviving edges would have required
$a$ to be present already.  Edges with zero modules impose no condition.
Consequently the chosen value at $a$ is compatible with every previous
value.  Moreover, the enlarged set is still upper closed under surviving
edges, since all vertices above $a$ are present.
The finite induction extends the original section to $[1,b]$.
Restriction to $K$ proves the required surjectivity.
\end{proof}

For $c\in K$, write $\mathcal O_{c,R}$ for the model in
Lemma~\ref{loc:models} with coefficients localized to $R$, restricted to
$K$.  It is supported on the full relative lower interval $[m,c]_D$,
which lies in $K$.  These models retain generation and relative
flabbiness.  To see the latter directly, a relative upper subset of
$K$ has the same intersection with $[m,c]_D$ as its upward closure in
the full coset.  A section on that subset extends by zero to this upward
closure outside the model's support, and then extends by the flabbiness
in Lemma~\ref{loc:models}.  Restrict the resulting section to $K$.
Localization commutes with these finite-support section conditions.

\subsection{The decomposition over the local ring}

Braden and MacPherson decompose pure sheaves by descending through the
vertices and splitting each free stalk over the image of sections above it
\cite[Theorem~1.4]{BM01}.  The next proposition carries out this descent
over the ungraded local ring $R$, retaining both endpoint restrictions on
every edge.

\begin{proposition}[Plane localization]
\label{loc:decomposition}
For each coset $C$ with $K=C\cap[1,b]\ne\varnothing$, there are
nonnegative integers $m_c$, indexed by $c\in K$, such that
\[
 B_R|_K\cong
 \bigoplus_{c\in K}\mathcal O_{c,R}^{\oplus m_c}
\]
as ungraded sheaves over $R$.
\end{proposition}

\begin{proof}
We construct an isomorphism in descending relative order.  Suppose the
sheaf has already been identified above a vertex $a$ with a sum of the
models in the statement.  More precisely, take any descending linear
extension of the finite poset $K$; the isomorphism has been constructed
on all previously processed vertices and on edges between them.
All relative upper neighbors of $a$ have already been processed.

Let $E_a$ be the set of relative upward edges from $a$ in $K$ and let
$K_{>a}=\{y\in K:a<_D y\}$.  Define the boundary image
\[
 N_a=\im\left(
    \Gamma(K_{>a},B_R)
       \longrightarrow\bigoplus_{e\in E_a}B_R^e
             \right),
\]
where the map restricts a section from the upper endpoint of each edge.
The image of the stalk map
\[
 \varphi:B_R^a\longrightarrow\bigoplus_{e\in E_a}B_R^e
\]
is exactly $N_a$.
For one inclusion, generation at $a$ extends any stalk value to a
section, whose upper restriction gives its boundary value.
For the other, relative flabbiness extends any section on $K_{>a}$ to
$K$, producing a preimage of its boundary value at $a$.
Thus $\varphi$ is a surjection from the free stalk $B_R^a$ onto $N_a$.

We will show that the stalks of the models already constructed above $a$
give a free cover of $N_a$ which is an isomorphism modulo $\mathfrak m$.
Splitting the given stalk through this cover will leave a free summand
with zero boundary map.  That summand will supply the models with top $a$.

The edge modules immediately above $a$, together with their upper
restriction maps, are already identified: the upper-end quotient rule
identifies each with its upper stalk modulo the corresponding label.
Hence the isomorphism already constructed above $a$ also identifies
$N_a$ with the boundary image for the sum of the existing models.
Only models with top $c>_D a$ contribute to this image.  Models whose
tops are not above $a$ have neither a stalk at $a$ nor a nonzero edge
module immediately above it.

Consider one contributing copy with top $c$.  Let $E_{a,c}$ be its
upward edges at $a$ whose upper endpoints are at most $c$ in relative
order.  This set is nonempty.  Its stalk map is
\[
 R\longrightarrow\bigoplus_{e\in E_{a,c}}R/(\alpha_e).
\]
Distinct labels at a vertex are nonassociate irreducibles in $R$.
Consequently its kernel is
\begin{equation}
\label{loc:boundary-product}
 \bigcap_{e\in E_{a,c}}\alpha_eR
  =p_{a,c}R,
 \qquad
 p_{a,c}=\prod_{e\in E_{a,c}}\alpha_e.
\end{equation}
The boundary image of this model is therefore $R/(p_{a,c})$.
Here we use its generation and flabbiness, as above, to identify the
stalk image with the boundary image coming from sections above $a$.
Since all the labels lie in $\mathfrak m$ and $E_{a,c}$ is nonempty,
$p_{a,c}\in\mathfrak m$.
Thus
\[
 R\longrightarrow R/(p_{a,c})
\]
is a minimal free cover: reduction modulo $\mathfrak m$ is an
isomorphism.

The section and boundary maps of a finite direct sum split componentwise.
Taking the sum of the existing model stalks at $a$ therefore gives a free
module $P_a$ and a surjection
\[
 \pi:P_a\longrightarrow N_a
\]
whose reduction modulo $\mathfrak m$ is an isomorphism.
If no model contributes, both $P_a$ and $N_a$ are zero and the same
argument applies.

We now split the given stalk by an explicit lift.
Because $B_R^a$ is free and $\pi$ is surjective, there exists
\[
 g:B_R^a\longrightarrow P_a,
 \qquad \pi g=\varphi.
\]
Modulo $\mathfrak m$, the map $\pi$ is an isomorphism and
$\varphi$ is surjective, so $g$ is surjective modulo $\mathfrak m$.
Nakayama's lemma applied to its finitely generated cokernel implies
that $g$ is surjective.  Since $P_a$ is free, $g$ splits.  We obtain
\[
 B_R^a\cong P_a\oplus Q_a,
 \qquad
 \varphi=(\pi,0),
\]
where $Q_a=\ker g$ is finite projective and therefore free over the
local ring $R$.

The factor $P_a$ extends the isomorphism for the existing models to
the stalk at $a$ and all its upward restrictions.  Add one copy of
$\mathcal O_{a,R}$ for each basis element of $Q_a$.  These new models
have zero outgoing edge maps at their top and vanish at every previously
processed vertex.  Thus they realize the factor $Q_a$ without changing
any previously constructed data.  The incoming edge modules will be
identified when their lower endpoints are processed, using the
upper-end quotient rule.
This completes the descending induction and proves the decomposition.
\end{proof}

\begin{remark}
\label{loc:arbitrary-edge-subsets}
The decomposition concerns the sheaf itself and does not depend on any
ordering of its edges.  In a summand $\mathcal O_{c,R}$ with $a\le_D c$,
the kernel conditions for any subset $F$ of upward edges at $a$ impose
exactly
\[
 L_a(F)=\left(
       \prod_{\substack{e\in F\\
                 \text{upper endpoint of }e\ \le_D c}}
                    \alpha_e\right)R,
\]
with the empty product equal to one.  Edges outside the support have
zero modules and impose no condition.  This identity is simply the
unique-factorization calculation in \eqref{loc:boundary-product}; in
particular it is valid for an edge order transported from another
Coxeter system.
\end{remark}
